% page 180 of the facsimile (image p-179 of the scan)
% block 157.5 x 242.0 mm ; left margin 8.0 mm
\pgc{-12.700mm}{0.000mm}{
\l{\cel{2}170.}
\l{}
\l{}
\l{\sou{flundro}. (\sou{zool.})\cel{2}Nomo di fisho de la genero \textquotedbl{}pleuronektoidi\textquotedbl{}.}
\l{}
\l{\sou{fluoro}. (\sou{kemio}).Korpo simpla, metaloida, qua deskompozas aquo}
\l{\cel{3}ye la temperaturo ordinara, e korodas preske omna metali.}
\l{\cel{3}-\cel{1}\sou{Simbolo kemiala}\cel{1}:\cel{2}F.\cel{3}- DEFIRS.}
\l{}
\l{\sou{fluorecar}. (\sou{netrans.}) (\sou{pri ula dissolvuri)}. Kovreskar ye en-}
\l{\cel{3}velopilo lumoza qua aspektas analoge kam opalo kande onu}
\l{\cel{3}lumizas li per suno-radii en chambro obskura. - DEFIRS.}
\l{}
\l{\sou{fluorino}. (\sou{kemio).}\cel{2}Kombinuro de fluoro e kalio. -\cel{1}\sou{Simbolo}}
\l{\cel{3}\sou{kemiala}\cel{1}: Ca F}
\l{}
\l{\sou{fluro}. (\sou{arkitekt.})\cel{2}Platformo en eskalero, ye la punto ube}
\l{\cel{3}finas etajo. - DE.}
\l{}
\l{\sou{fluto}. (\sou{muziko)}. Muzik-instrumento ventala, cilindro apertata}
\l{\cel{3}supre, klozata infre, qua havas, apud la parto supra,}
\l{\cel{3}aperturo qua uzesas kom boko-peco ube onu suflas, e longe-}
\l{\cel{3}sale provizita ye trui sur qui onu pozas la fingri por}
\l{\cel{3}apertar e klozar ta trui segun la noto quan onu volas pro-}
\l{\cel{3}duktar. - DEFIRS.}
\l{}
\l{\sou{fluvio}. I. Granda aquofluo qua konservas sua nomo til la maro}
\l{\cel{3}ube lu debushas. - II. (\sou{metaf}.)To quo fluas abunde quale}
\l{\cel{3}fluvio. - EFIRS.}
\l{}
\l{\sou{fluxo}. La mareo adlitora. - EFIRS.}
\l{}
\l{\sou{fluxiono.}\cel{2}I. (\sou{patol.}) Adfluo di sango o di altra liquido ad}
\l{\cel{3}ula tisui, ad ula organi, konseque de stando inflamala. -}
\l{\cel{3}II. (\sou{matem.}) La quanti finita quin produktas movo o fluo}
\l{\cel{3}kontinua (la lineo, per la movo di la punto : la surfaco,}
\l{\cel{3}per la movo di la lineo, e c.) - DEFIS.}
\l{}
\l{\sou{foko}.\cel{2}I. (\sou{optiko}).Punto a qua konvergas la lumo-radii reflek-}
\l{\cel{3}tita o refraktita. - II. (\sou{geom}.)Punto o punti de ube de-}
\l{\cel{3}partas ed adube su asemblas la vektori. - DEFIR.}
\l{}
\l{\sou{fola}. Termino vulgara qua aplikesas precipue a la formi maxim}
\l{\cel{3}videbla e shokiva di alienaco, a ti qui su manifestas per}
\l{\cel{3}parolo sencesa, agiteso e violento. - EFI.}
\l{}
\l{\sou{folado}. (\sou{zool.)}\cel{2}Molusko bivalva qua perforas la roki. - L.}
\l{\cel{3}\sou{pholadidea}. -}
\l{}
\l{\sou{folio}. I. (\sou{bot.}) Lamelo verda qua naskas de la stipi, de la}
\l{\cel{3}rami. II.To quo esas dina e plata quale arboro-folio}
\l{\cel{3}(papero, e c.) - EFIRS.}
\l{}
\l{\sou{folietono.}\cel{3}Parto rezervita, infre de jurnalo-pagino (cirkume}
\l{\cel{3}la triimo di la alteso) por la publikigo di artikli kriti-}
\l{\cel{3}kala, di romani, e c. - DEFRS.}
}
